% !TeX root = ../main.tex
% =========================================================================
% 7. METODOLOGÍA / PROPUESTA DE DESARROLLO
% Párrafo 1: taxonomía investigativa. Fases continuas.
% Protocolo experimental cuantitativo formal (métricas, hipótesis, muestra).
% Figura en TikZ con caption y nota APA 7.
% =========================================================================
\section{Metodología / Propuesta de desarrollo}

La presente investigación adopta un alcance descriptivo y aplicado, bajo un enfoque metodológico mixto predominantemente cuantitativo y con un diseño cuasi-experimental de corte transversal. Según su propósito, constituye una investigación tecnológica orientada a la concepción, ingeniería y validación empírica de un artefacto de software; según el método de inferencia, es de naturaleza deductiva al derivar hipótesis operativas a partir de constructos teóricos consolidados; según la procedencia de la información, combina fuentes secundarias documentales (catálogos de fabricantes, especificaciones técnicas y literatura científica indexada) con datos primarios generados durante las pruebas automatizadas del sistema. La recolección de datos se sustentará en el registro automatizado de trazas de telemetría y métricas de rendimiento en servidor, complementado con entrevistas semiestructuradas al equipo comercial para el levantamiento de requerimientos de dominio \parencite{Llerena-Izquierdo_Ayala-Carabajo_2022}.

El desarrollo del proyecto se estructurará rigurosamente en cinco fases consecutivas e iterativas:

En la \textbf{Fase 1: Revisión bibliográfica y análisis de requerimientos}, se llevará a cabo una revisión bibliográfica exhaustiva fundamentada en directrices metodológicas para ingeniería de software y la recopilación de literatura científica reciente en bases de datos indexadas (Scopus, IEEE Xplore, ACM y Web of Science). Se analizarán investigaciones sobre agentes conversacionales, modelos de lenguaje extenso y técnicas de recuperación aumentada (RAG). De forma paralela, se levantarán los requerimientos funcionales y técnicos con el personal comercial de la distribuidora Power Electrimk, estructurando un catálogo industrial multimarca inicial superior a 600 referencias activas y verificadas de fabricantes líderes (Schneider Electric y Siemens) distribuidas en seis líneas esenciales: distribución de potencia y modular, control y maniobra de motores, automatización programable (PLCs y módulos lógicos), variadores de frecuencia, fuentes de poder y unidades de mando/señalización.

En la \textbf{Fase 2: Diseño de la arquitectura propuesta}, se definirá la arquitectura de software en capas basada en microservicios ligeros, estableciendo el esquema relacional en SQLite optimizado bajo modo WAL y dotado de una capa de búsqueda de texto completo (\textit{Full-Text Search}, FTS5) con algoritmo de clasificación BM25 y rutinas de auto-reparación dinámica de índices. Se diseñará el flujo de control para la verificación y consumo de los webhooks de WhatsApp Cloud API, la arquitectura de defensa en profundidad contra inyecciones de \textit{prompt} en documentos técnicos e imágenes fotográficas, el mecanismo de \textit{tokens} centinela efímeros (\textit{canary tokens}) para certificar la integridad de las directivas de sistema, y la política comercial multimarca 'solución primero, marca bajo demanda' orientada a prevenir la fatiga de decisión en el cliente \parencite{Zerega-Prado2022}.

En la \textbf{Fase 3: Desarrollo del prototipo funcional}, se programará el servicio web en lenguaje C\# sobre la plataforma .NET 8, aprovechando sus características nativas de alto rendimiento asíncrono y bajo consumo de memoria. Se implementará la lógica del embudo comercial (calificación de prospectos, asesoría técnica de compatibilidad, consulta determinista de inventario y derivación formal hacia asesores comerciales humanos sin alucinación de precios), el módulo de persistencia y contexto conversacional por número de teléfono, el pipeline de inspección segura en memoria RAM (Zero Disk) para documentos de requerimientos de proyectos (.xlsx, .csv, .pdf) con neutralización de inyección de fórmulas y entidades externas XML (XXE), el procesamiento multimodal nativo de fotografías técnicas para identificación de placas de motores y referencias legadas, el conector HTTP con autenticación Bearer hacia la API de WhatsApp, la capa de consulta relacional SQLite WAL con FTS5 y la integración con el SDK de Google Gemini para la síntesis contextual en lenguaje natural \parencite{educsci14121330}. La arquitectura se estructurará mediante interfaces desacopladas que habilitan contractualmente la integración de seguimiento automatizado de cotizaciones y sincronización con plataformas CRM en fases operativas posteriores.

En la \textbf{Fase 4: Despliegue y pruebas experimentales}, el prototipo se desplegará en una instancia virtual Linux (Ubuntu 24.04.5 LTS) en Amazon Web Services (AWS EC2 t2.micro, acogida a la Capa Gratuita de 750 horas mensuales), configurando un servidor Nginx como proxy inverso hacia el runtime Kestrel en puerto 5252, administrado por el demonio Systemd (\texttt{agenteventas.service}) y con terminación de certificados de seguridad TLS 1.3 emitidos por Certbot (Let's Encrypt) bajo el subdominio oficial \texttt{agente.powerelectrimk.com}. Se ejecutará un protocolo experimental automatizado compuesto por una batería de $N = 50$ casos de prueba estructurados (búsquedas por código exacto, por especificación técnica, por categoría y solicitudes multimarca), comparando el desempeño del agente con recuperación local determinista frente al modelo generativo sin recuperación para cuantificar la tasa de eliminación de alucinaciones en precios y referencias de catálogo.

En la \textbf{Fase 5: Análisis de resultados y documentación}, se procesarán estadísticamente las métricas de latencia de respuesta, porcentaje de exactitud en precios y referencias técnicas, tasa de alucinación y utilización de memoria RAM y procesador en el servidor. Los resultados se contrastarán contra las hipótesis de investigación planteadas, y se redactará el informe técnico final de titulación de conformidad con las directrices institucionales \parencite{10.1007/978-3-031-98890-5_3}.

La Figura \ref{fig:metodologia} sintetiza el flujo de fases del proyecto y los entregables correspondientes a cada etapa de la investigación.

\begin{figure}[H]
    \centering
    \caption{Diagrama de flujo de las fases de la metodología de investigación}
    \label{fig:metodologia}
    \begin{tikzpicture}[
        node distance=0.55cm,
        fase/.style={rectangle, rounded corners=3pt, draw=blue!60!black, fill=blue!8,
                     text width=4.4cm, minimum height=1.45cm, align=center, font=\small},
        prod/.style={rectangle, draw=gray!70, fill=gray!8, text width=5.6cm,
                     minimum height=1.45cm, align=left, font=\footnotesize},
        flecha/.style={-{Stealth[length=2.5mm]}, thick, blue!60!black}
    ]
        \node[fase] (f1) {\textbf{Fase 1}\\Revisión bibliográfica y análisis de requerimientos};
        \node[fase, below=of f1] (f2) {\textbf{Fase 2}\\Diseño de la arquitectura propuesta};
        \node[fase, below=of f2] (f3) {\textbf{Fase 3}\\Desarrollo del prototipo funcional};
        \node[fase, below=of f3] (f4) {\textbf{Fase 4}\\Despliegue y pruebas experimentales};
        \node[fase, below=of f4] (f5) {\textbf{Fase 5}\\Análisis de resultados y documentación};

        \node[prod, right=1.2cm of f1] (p1) {Estado del arte científico, matriz de requerimientos y catálogos normalizados.};
        \node[prod, right=1.2cm of f2] (p2) {Diseño arquitectónico, modelo relacional SQLite WAL y especificación de webhooks.};
        \node[prod, right=1.2cm of f3] (p3) {Servicio .NET 8 con Gemini, recuperación local SQLite y memoria de sesión.};
        \node[prod, right=1.2cm of f4] (p4) {Instancia AWS con Nginx y TLS 1.3; batería de 50 pruebas y telemetría registrada.};
        \node[prod, right=1.2cm of f5] (p5) {Evaluación comparativa (con/sin RAG), reporte de hipótesis e informe final.};

        \foreach \a/\b in {f1/f2, f2/f3, f3/f4, f4/f5} \draw[flecha] (\a) -- (\b);
        \foreach \i in {1,...,5} \draw[-{Stealth[length=2mm]}, gray!70] (f\i) -- (p\i);
        \draw[flecha, dashed] (f4.west) -- ++(-0.6,0) |- node[pos=0.25, left, font=\scriptsize, text width=1.8cm, align=center] {ajustes y\\calibración} (f3.west);
    \end{tikzpicture}
    \notatabla{El diagrama ilustra la progresión secuencial de las cinco fases metodológicas y sus entregables tangibles, incluyendo el lazo de retroalimentación de ajustes entre la fase experimental y el código del prototipo.}
\end{figure}

Para la evaluación formal del sistema se han operacionalizado las siguientes métricas experimentales:
\begin{itemize}
    \item \textbf{Latencia de respuesta ($T_{\text{resp}}$):} Tiempo total transcurrido en segundos desde la recepción del evento HTTP en el webhook hasta la emisión de la llamada de respuesta a la API de WhatsApp, medido con precisión de milisegundos mediante marcas de tiempo en el servidor.
    \item \textbf{Exactitud de cotización ($Acc_{\text{cot}}$):} Porcentaje de cotizaciones emitidas donde los códigos de producto, marcas y precios unitarios coinciden exactamente con los registros vigentes del inventario oficial, definido como:
    \begin{equation}
        Acc_{\text{cot}} = \left( \frac{C_{\text{correctas}}}{C_{\text{totales}}} \right) \times 100\%
    \end{equation}
    \item \textbf{Tasa de alucinación ($H_{\text{rate}}$):} Proporción de respuestas en las que el modelo generativo asevera precios inexistentes, inventa atributos de productos o recomienda ítems no catalogados en la base de datos comercial.
    \item \textbf{Consumo de recursos:} Perfil de utilización de memoria RAM del proceso .NET (en megabytes) y porcentaje promedio de ocupación de CPU en la instancia de nube durante ráfagas de consultas.
\end{itemize}
